\DocumentMetadata{lang=es-DO,testphase={phase-III,math}}
\documentclass[11pt]{scrartcl}
\usepackage{icv}

% -----------------------------------------------------------------------
% Metadata -- reemplaza estos valores para cada handout nuevo.
% -----------------------------------------------------------------------
\icvsetup{
  asignatura  = {ICV 442-T -- Acueductos y Alcantarillados},
  titulocorto = {Acueductos y Alcantarillados},
  titulo      = {Título del handout},
  subtitulo   = {Subtítulo o tema específico (opcional)},
  profesor    = {Dr. Ricardo Hernández Moreira},
  periodo     = {2026-1},
  version     = {v1},
}

\begin{document}
\icvmaketitle

\begin{icvobjectives}
  \item Primer objetivo de aprendizaje de este handout.
  \item Segundo objetivo de aprendizaje.
\end{icvobjectives}

\section{Introducción}

Contexto breve: qué problema resuelve este handout y cuándo conviene
consultarlo.

\section{Procedimiento}

Pasos o desarrollo del método, con ecuaciones cuando aplique:
\[
  Q = A \cdot V
\]
donde \(Q\) es el caudal, \(A\) el área de la sección transversal y \(V\)
la velocidad media.

\section{Ejemplo resuelto}

Ejemplo numérico con unidades vía \texttt{siunitx}: un caudal de
\qty{150}{\liter\per\second} equivale a \qty{0.15}{\cubic\meter\per\second}.

\begin{table}[htbp]
  \centering
  \caption{Ejemplo de tabla con columnas alineadas al decimal}
  \begin{tabular}{l S[table-format=3.1]}
    \toprule
    {Sección} & {Caudal (\unit{\liter\per\second})} \\
    \midrule
    A & 120.5 \\
    B & 98.3  \\
    C & 150.0 \\
    \bottomrule
  \end{tabular}
\end{table}

\section{Observaciones de uso}

Notas prácticas, casos límite o errores comunes al aplicar este método.

\end{document}
